% -*- coding: utf-8 -*-
% !TEX program = xelatex
\documentclass[plain,noanswer,medmath]{../../template/jnuexam}
\usepackage{../../template/kysx}

\begin{document}


\examtitle{1988}{3}


\loadproblems{1}{1988P1}%载入试卷一


\makepart{填空题}{本题满分20分，每小题4分}


\begin{problem}
若$f(x)=\begin{cases}
  \e^2(\sin x+\cos x), & x>0 \\
                 2x+a, & x\le 0
\end{cases}$是$(-\infty,+\infty)$上的连续函数，则$a=$ \fillin{}．
\end{problem}


\useproblem{1}{2}{1}%【同试卷一第二(1)题】


\useproblem{1}{2}{3}%【同试卷一第二(3)题】


\begin{problem}
$\lim_{x\to +0}\,\left(\frac{1}{\sqrt{x}} \right)^{\tan x}=$ \fillin{}．
\end{problem}


\begin{problem}
$\int_0^4\e^{\sqrt{x}} \dx=$ \fillin{}．
\end{problem}


\makepart{选择题}{本题满分20分，每小题4分}


\begin{problem}
$f(x)=\frac{1}{3}x^3+\frac{1}{2}x^2+6x+1$的图形在点$\left(0,1 \right)$处切线与$x$轴交点的坐标是\pickout{}
\begin{abcd}
\item $\left(-\frac{1}{6},0\right)$．         
\item $\left(-1,0\right)$．      
\item $\left(\frac{1}{6},0\right)$．      
\item $\left(1,0\right)$．
\end{abcd}
\end{problem}


\begin{problem}
若$f(x)$与$g(x)$在$(-\infty,+\infty)$上皆可导，且$f(x)<g(x)$，则必有\pickout{}
\begin{abcd}
\item $f(-x)>g(-x)$．
\item $f'(x)<g'(x)$．
\item $\lim_{x\to x_0}f(x)<\lim_{x\to x_0}g(x)$．
\item $\int_0^x f(t)\dt<\int_0^x g(t)\dt$．
\end{abcd}
\end{problem}


\useproblem{1}{3}{1}%【同试卷一第三(1)题】


\begin{problem}
曲线$y=\sin^{3/2} x$（$0\le x\le \pi$）
与$x$轴围成的图形绕$x$轴旋转所形成的旋转体体积为\pickout{}
\begin{abcd}
\item $\frac{4}{3}$．          
\item $\frac{4}{3}\pi$．         
\item $\frac{2}{3}\pi^2$．      
\item $\frac{2}{3}\pi$．
\end{abcd}
\end{problem}


\useproblem{1}{3}{5}%【同试卷一第三(5)题】


\makepart{计算题}{本题满分15分，每小题5分}


\useproblem{1}{1}{2}%【同试卷一第一(2)题】


\begin{problem}
已知$y=1+x\e^{xy}$，求$y'|_{x=0}$及$y''|_{x=0}$．
\end{problem}


\begin{problem}
求微分方程$y'+\frac{1}{x}y=\frac{1}{x(x^2+1)}$的通解（一般解）．
\end{problem}


\makepart{}{本题满分12分}

\begin{problem}
作函数$y=\frac{6}{x^2-2x+4}$的图形，并填写下表．
\par\centerline{\begin{tabularx}{0.9\linewidth}{|l|X|}
\hline
 单调增区间  &   \\  
\hline
 单调减区间  &   \\  
\hline
 极值点  &   \\
\hline
 极值  &   \\
\hline
 凹$(\cup)$区间  &   \\
\hline
 凸$(\cap)$区间  &   \\
\hline
 拐点 &   \\
\hline
 渐近线  &   \\
\hline
\end{tabularx}}
\end{problem}


\makepart{}{本题满分8分}

\begin{problem}
将长为$a$的铁丝切成两段，一段围成正方形，另一段围成圆形．问这两段铁丝各长为多少时，正方形与圆形的面积之和为最小？
\end{problem}


\makepart{}{本题满分10分}%【同试卷一第五题】

\useproblem{1}{5}{0}


\makepart{}{本题满分6分}

\begin{problem}
设$x\ge -1$，求$\int_{-1}^x \left(1-|t| \right)\dt$．
\end{problem}


\makepart{}{本题满分6分}

\begin{problem}
设$f(x)$在$(-\infty +\infty)$上有连续导数，且$m\le f(x)\le M$．
\begin{enumerate}
\item 求$\lim_{a\to +0}\,\frac{1}{4a^2}\int_{-a}^a [f(t+a)-f(t-a)]\dt$．
\item 证明$\left| \frac{1}{2a}\int_{-a}^a f(t)\dt-f(x) \right|\le M-m$（$a>0$）．
\end{enumerate}
\end{problem}


\end{document}
